% Chapter 6: Maryam Mirzakhani's Work on Hyperbolic Surfaces
\let\LocalMacro\relax

\chapter{Maryam Mirzakhani's Work on Hyperbolic Surfaces}

\section{The Problem}

\subsection{Informal}
Some surfaces curve inward everywhere like the inside of a ball, and some curve outward everywhere like the outside of a sphere. But there is a third kind that curves inward in one direction and outward in another, like a saddle or a Pringle chip. These saddle-shaped surfaces come in infinitely many variations, and the mathematician Maryam Mirzakhani wanted to understand the complete landscape of all possible versions. She figured out how to measure the size of this landscape and discovered surprising formulas connecting loops drawn on the surface to its overall shape.

\subsection{Formal}

The central question is: what are the volumes of moduli spaces of hyperbolic surfaces with respect to the Weil--Petersson metric? More precisely, for a surface of genus $g$ with $n$ geodesic boundary components of lengths $L_1, \ldots, L_n$, is the volume $V_{g,n}(L_1, \ldots, L_n)$ a polynomial in the $L_i$? And can one find a recursive formula to compute it?

\subsection{Historical Context}
Mirzakhani was the first woman to win the Fields Medal (2014). She earned two gold medals at the International Mathematical Olympiad (1994, 1995) before becoming a professor at Stanford University. Her work on hyperbolic surfaces was recognized as among the most beautiful and profound in modern geometry. She passed away in 2017 at age 40, leaving an extraordinary legacy.

\subsection{What Was Known}
Before Mirzakhani, the volumes of moduli spaces were known only for a few small cases. Mathematicians knew the Weil--Petersson metric was well-defined and had some basic properties, but there was no general formula for the volumes. The growth of simple closed geodesics was known to be polynomial, but the exact growth rate and its connection to the moduli space were unknown.

\subsection{Why It Matters}
Hyperbolic geometry appears throughout modern mathematics and physics, from the study of the universe's shape to the behavior of random networks. Mirzakhani's formulas for the volumes of moduli spaces give mathematicians a way to measure and compare these infinite landscapes of shapes, connecting geometry to algebra and physics in unexpected ways.

\subsection{Simplified Version}
The simplest case is a once-punctured torus (a torus with one small hole), which has genus 1 and one boundary component. Its moduli space volume was already computed by Mirzakhani herself, and the polynomial structure of these volumes was verified for small genera. The full conjecture asks whether this polynomial structure persists for all genera and all boundary configurations.

\section{Mathematical Theory}


\subsection{Prerequisites}
This chapter assumes familiarity with:
\begin{itemize}
\item Differential geometry (surfaces, curvature, geodesics)
\item Teichm\"uller theory (moduli spaces, mapping class groups)
\item Hyperbolic geometry (constant negative curvature, $\mathbb{H}^2$)
\item Basic algebraic topology (Euler characteristic, intersection theory)
\end{itemize}

\subsection{The Moduli Space and Weil--Petersson Metric}

The moduli space $\mathcal{M}_{g,n}$ of hyperbolic surfaces of genus $g$ with $n$ punctures (or geodesic boundaries of lengths $L_1, \ldots, L_n$) carries a natural Riemannian metric called the Weil--Petersson metric.

\begin{definition}[Weil--Petersson Metric]
The \emph{Weil--Petersson metric} on Teichm\"uller space $\mathcal{T}_{g,n}$ is the K\"ahler metric defined by
\[\langle \mu, \nu \rangle_{WP} = \int_{\Sigma} \mu \overline{\nu} \, \rho^2 \, dA\]
where $\mu, \nu$ are Beltrami differentials representing tangent vectors, and $\rho^2|dz|^2$ is the hyperbolic metric on the surface $\Sigma$.
\end{definition}

This metric is incomplete: surfaces can degenerate by pinching geodesics, and the distance to the boundary is finite.

\subsection{Simple Closed Geodesics}

A simple closed geodesic on a hyperbolic surface is a closed curve that locally minimizes distance and does not intersect itself. The set of all simple closed geodesics up to isotopy is countable.

Mirzakhani showed that the number of simple closed geodesics of length at most $L$ grows polynomially in $L$:

\begin{equation}
s_\gamma(L) \sim c_\gamma \cdot L^{6g-6}
\end{equation}

where $c_\gamma$ depends on the topology of the curve $\gamma$ (its position in the space of measured laminations).

\subsection{Mirror Symmetry and Intersection Theory}

Mirzakhani's breakthrough was a deep connection between counting geodesics and intersection theory on the moduli space. She expressed the growth constants $c_\gamma$ in terms of integrals of tautological classes on $\mathcal{M}_{g,n}$.

The \emph{tautological classes} are cohomology classes $\kappa_i, \psi_i, \lambda_i$ on $\mathcal{M}_{g,n}$ that encode geometric information. For example:
\begin{itemize}
\item $\psi_i = c_1(\mathbb{L}_i)$ where $\mathbb{L}_i$ is the cotangent line bundle at the $i$-th marked point.
\item $\kappa_i$ are the Miller--Morita--Mumford classes, pushforwards of $\psi$-classes.
\end{itemize}

\section{The Solution}

\subsection{The Big Idea}
Imagine a vast landscape of possible saddle-shaped surfaces, like an infinite collection of Pringle chips with different curvatures. Mirzakhani wanted to measure the total ``size'' of this landscape---not just compute it for one surface at a time, but find a formula that works for \emph{all} of them simultaneously. Her key metaphor: imagine gently shearing the surface along its internal curves, like sliding a deck of cards sideways. This shearing motion is called an earthquake flow. By studying how the earthquake flow rearranges the landscape, she discovered that the volume of complicated surfaces can always be built up from the volumes of simpler ones---like assembling a mosaic from smaller tiles. The result is a master equation: a polynomial recursion that works for every genus at once.

\subsection{What Was Stopping Progress}
Before Mirzakhani, mathematicians had managed to compute the ``size'' (volume) of the landscape for just a handful of the simplest surfaces---like measuring the area of only the first few rooms in an infinite building. There was no general formula, not even a guess, for what the volume should look like for surfaces of arbitrary complexity. Existing tools from Riemannian geometry could give rough estimates of volume but not exact answers, and tools from algebraic geometry could count specific objects on surfaces but didn't connect to the actual geometric volume. It was as if two teams of surveyors were measuring different properties of the same territory but had no map linking their measurements together.

\subsection{The Breakthrough}
Mirzakhani connected three worlds that nobody had linked before: hyperbolic geometry (counting curves on surfaces), differential geometry (measuring volumes of shape-spaces), and algebraic geometry (intersection numbers that arise in string theory). She did this by exploiting the earthquake flow---a natural dynamical system on Teichmuller space---to decompose the volume integral into pieces that correspond to ``pinching'' geodesics. Each piece factors into contributions from lower-genus surfaces, yielding a recursion. The deep surprise: the volume is always a \emph{polynomial} in the boundary lengths, and its coefficients encode mirror-symmetry data.

\subsection{How It Works}

\subsection{What Was Missing}
Before Mirzakhani, the volumes $V_{g,n}(L_1, \ldots, L_n)$ of moduli spaces of bordered Riemann surfaces with respect to the Weil--Petersson metric were known only for a handful of low-genus cases. There was no general formula---not even a conjectured form---for arbitrary $g$ and $n$. The fundamental question was: do these volumes follow a predictable pattern (such as being polynomials in the $L_i$), and if so, can they be computed systematically? The gap was total: no recursion, no closed-form, and no structural theorem existed.

\subsection{Why Existing Methods Failed}
Prior approaches relied on topological recursion and matrix models, which gave \emph{recursive structures} for certain enumerative invariants (such as intersection numbers on $\overline{\mathcal{M}}_{g,n}$) but did not directly yield explicit polynomial formulas for Weil--Petersson volumes. The problem was that these methods operated at the level of cohomological invariants without connecting to the \emph{dynamics} of the moduli space. One could count geodesics or compute intersection numbers, but the link between the geometry of Teichm\"uller space and the combinatorics of the volume polynomial was missing. Existing techniques from Riemannian geometry (comparison theorems, heat kernels) gave asymptotic estimates but not exact formulas.

\subsection{What Was Novel}
Mirzakhani's key innovation was the \emph{earthquake flow}: she used the dynamics of earthquakes (continuous left-right shearing along measured geodesic laminations) on Teichm\"uller space to derive a recursion for Weil--Petersson volumes. The earthquake flow preserves the WP symplectic form and acts ergodically, allowing Mirzakhani to express the volume as an integral over the space of measured laminations. By decomposing this integral according to how geodesics are pinched, she obtained a recursion that expresses $V_{g,n}$ in terms of volumes of lower-genus surfaces. This connected three seemingly unrelated worlds: hyperbolic geometry (counting geodesics), differential geometry (WP volumes), and algebraic geometry (intersection numbers and mirror symmetry).

\subsection{Student-Friendly Explanation}
Imagine you have a vast landscape of possible saddle-shaped surfaces, each described by how its curves are arranged. Mirzakhani asked: can we measure the total ``size'' of this landscape? Previous mathematicians had computed this size for a few tiny pieces but had no recipe for the whole thing. Mirzakhani's idea was to use a natural motion on the landscape---imagine gently shifting the surface by sliding it along its curves, like shearing a deck of cards. By studying how this motion redistributes the landscape, she found a master equation (a recursion) that builds the volume for complicated surfaces out of volumes of simpler ones. The result is a beautiful polynomial formula, and the method connects hyperbolic geometry to mirror symmetry in a way nobody had expected.

\subsection{Recursive Formula for Weil--Petersson Volumes (2007--2008)}

\begin{theorem}[Mirzakhani, 2007 \cite{mirzakhani2007}]
The volume of the moduli space $\mathcal{M}_{g,n}(L_1, \ldots, L_n)$ with respect to the Weil--Petersson metric satisfies a recursive formula:
\begin{equation}
V_{g,n}(L_1, \ldots, L_n) = \int_{\mathcal{M}_{g,n}(L_1, \ldots, L_n)} e^{-\omega_{WP}}
\end{equation}
where $V_{g,n}$ is a polynomial in the boundary lengths $L_i$ of degree $3g-3+n$.
\end{theorem}

The recursion is:
\begin{multline}
V_{g,n}(L_1, \ldots, L_n) = \frac{1}{2} \sum_{i+j = n-1} \sum_{g_1+g_2 = g} \int_0^\infty \int_0^\infty V_{g_1,i+1}(L_1, \ldots, L_i, x, y) \\
\times V_{g_2,j+1}(L_{i+1}, \ldots, L_{i+j}, x, y) \, x y \, dx \, dy \\
+ \frac{1}{2} \int_0^\infty \int_0^\infty V_{g-1,n+1}(L_1, \ldots, L_n, x, y) \, x y \, dx \, dy
\end{multline}

This was revolutionary because it connected:
\begin{itemize}
\item Counting simple closed geodesics (combinatorial geometry)
\item Weil--Petersson volumes (differential geometry)
\item Mirror symmetry (algebraic geometry)
\end{itemize}

\begin{highlightbox}
Mirzakhani's recursive formula was unexpected because it connected the Weil--Petersson volumes to mirror symmetry in a way that had never been anticipated. Her methods combined techniques from hyperbolic geometry, Teichm\"uller theory, and algebraic geometry in a novel and powerful way.
\end{highlightbox}

\subsection{After the Proof}

Mirzakhani's results are completely settled. She proved the recursive formula for Weil--Petersson volumes, the polynomial structure of those volumes, the growth rate of simple closed geodesics, and a new proof of the Witten--Kontsevich theorem. Each of these results has been verified and extended by the mathematical community.

Her methods opened deep connections between hyperbolic geometry, mirror symmetry, and integrable systems. The recursive formula revealed that Weil--Petersson volumes encode data about mirror symmetry for curves, linking Teichm\"uller theory to string-theoretic physics in unexpected ways. Her counting results for geodesics have been extended to flat surfaces, translation surfaces, and higher-dimensional Teichm\"uller spaces, spawning active research programs.

Several important questions remain open. Explicit closed-form formulas for Weil--Petersson volumes in specific genera (beyond what the recursion produces) are still unknown. The deeper connections between Mirzakhani's work and quantum gravity---where moduli spaces of surfaces arise in path integrals---are still being explored. Generalizations to higher Teichm\"uller theory (for non-compact or higher-rank groups) and to surfaces with more complicated boundary conditions are active areas of ongoing research.

\subsection{Counting Simple Closed Geodesics (2008)}

Mirzakhani proved that the number of simple closed geodesics of length at most $L$ on a hyperbolic surface grows like $L^{6g-6}$, with the leading coefficient expressible as an intersection number on the moduli space \cite{mirzakhani2008}.

For a fixed simple closed curve $\gamma$ on a surface $\Sigma_{g,n}$:
\begin{equation}
\lim_{L \to \infty} \frac{s_\gamma(L)}{L^{6g-6}} = \frac{2^{|\gamma|} (6g-6)!}{\dim \mathcal{ML}_{g,n}} \int_{\mathcal{ML}_{g,n}} \cdots
\end{equation}

where $\mathcal{ML}_{g,n}$ is the space of measured laminations, and the integral involves Thurston's symplectic measure.

\subsection{The Witten--Kontsevich Theorem (2013)}

Mirzakhani gave a new proof of the Witten--Kontsevich theorem, which relates intersection numbers on $\mathcal{M}_{g,n}$ to the KdV hierarchy of integrable systems \cite{mirzakhani2013}. This established a bridge between hyperbolic geometry and integrable systems.

\subsection{Moments of Closed Geodesics (2014)}

In her Fields Medal lecture, Mirzakhani presented results on the moments of closed geodesics and their relationship to the Weil--Petersson volume polynomial. She showed that the $k$-th moment of the length of a random simple closed geodesic is a rational multiple of the Weil--Petersson volume.

\section{Impact}

\begin{itemize}
\item \textbf{Teichm\"uller theory}: Mirzakhani's recursive formulas for Weil--Petersson volumes have become fundamental tools. They allow computation of volumes to arbitrarily high genus.
\item \textbf{Counting geodesics}: Her methods for counting simple closed geodesics have been extended to many other settings, including counting curves on flat surfaces and on higher-dimensional varieties.
\item \textbf{Mirror symmetry}: The connection she discovered between counting geodesics and mirror symmetry opened new research directions linking Teichm\"uller theory, algebraic geometry, and mathematical physics.
\item \textbf{Ergodic theory}: Her work on the earthquake flow on moduli space has connections to Ratner's theorems and dynamics on homogeneous spaces.
\end{itemize}

\section{Exercises}

\begin{exercise}[Basic]
Show that the Euler characteristic of a surface of genus $g$ is $2 - 2g$. Why does this imply that surfaces of genus $g \geq 2$ admit hyperbolic metrics? \emph{Hint: Uniformization theorem.}
\end{exercise}

\begin{exercise}[Intermediate]
Explain why the Weil--Petersson metric is incomplete on Teichm\"uller space. What are the boundary points? \emph{Hint: Consider pinching a simple closed geodesic to length 0.}
\end{exercise}

\begin{exercise}[Challenge]
Derive the growth rate $L^{6g-6}$ for the number of simple closed geodesics on a surface of genus $g$. Why does the exponent depend only on the genus? \emph{Hint: The dimension of the space of measured laminations is $6g-6$.}
\end{exercise}


\begin{exercise}[Computational]
Write a Python/SageMath script to explore a concept from this chapter. See the appendix for starter code.
\end{exercise}

\section{Timeline}

\begin{tabularx}{\linewidth}{lX}
\toprule
\textbf{Year} & \textbf{Event} \\ \midrule
2004 & Mirzakhani's PhD thesis at Harvard (supervised by Curtis McMullen) \\ 
2007 & Mirzakhani proves the recursive formula for Weil--Petersson volumes \cite{mirzakhani2007} \\ 
2008 & Mirzakhani publishes counting geodesics results \cite{mirzakhani2008} \\ 
2013 & Mirzakhani proves the Witten conjecture/Kontsevich theorem \cite{mirzakhani2013} \\ 
2014 & Fields Medal awarded to Mirzakhani \\ 
2017 & Mirzakhani passes away at age 40 \\ \bottomrule
\end{tabularx}

\section{Citations}

\begin{enumerate}
\item Mirzakhani, M. (2007). ``Simple geodesics and Weil--Petersson volumes of moduli spaces of bordered Riemann surfaces.'' Inventiones Mathematicae, 167(1), 179--222.
\item Mirzakhani, M. (2008). ``Simple geodesics and Weil--Petersson volumes of moduli spaces of bordered Riemann surfaces.'' Inventiones Mathematicae, 171(1), 179--222.
\item Mirzakhani, M. (2013). ``A simple proof of the Witten--Kontsevich theorem.'' Geometry \& Topology, 17(2), 681--706.
\item Mirzakhani, M. (2014). ``Hyperbolic geometry on moduli space.'' Annals of Mathematics Studies, 183, Princeton University Press.
\item Wolpert, S. (2008). ``On the Weil--Petersson geometry of the moduli space of curves.'' American Journal of Mathematics, 130(4), 1045--1069.
\end{enumerate}